% Generated draft: requires mathematical and source review.
\documentclass[11pt]{article}
\usepackage[margin=25mm]{geometry}
\usepackage{fontspec}
\setmainfont{Latin Modern Roman}
\usepackage{amsmath,amssymb,mathtools,unicode-math}
\setmathfont{Latin Modern Math}
\usepackage{enumitem,needspace,xurl,hyperref}
\setlength{\parindent}{0pt}
\setlength{\parskip}{5pt}
\setlength{\emergencystretch}{3em}
\newcommand{\N}{\mathbb{N}}
\newcommand{\Z}{\mathbb{Z}}
\newcommand{\Q}{\mathbb{Q}}
\newcommand{\R}{\mathbb{R}}
\newcommand{\C}{\mathbb{C}}
\newcommand{\F}{\mathbb{F}}
\newcommand{\A}{\mathbb{A}}
\newcommand{\PP}{\mathbb{P}}
\newcommand{\E}{\mathbb{E}}
\newcommand{\eps}{\varepsilon}
\newcommand{\dd}{\,\mathrm{d}}
\begin{document}
\section*{UnsolvedMath ID: 30001003}
\textbf{Source problem code:} catalogue entry
\subsection{Definitions and mathematical statement}
Let \((X,L)\) be a polarized compact Kähler manifold, with \(L\) an ample holomorphic line bundle, and let \([\omega_0]\) be the Kähler class in the source. A representative \(\omega\in[\omega_0]\) has constant scalar curvature if \(\operatorname{Scal}(\omega)\) is constant on \(X\).

The source asks, as a conjecture attributed to Yau, whether the following equivalence holds:
\[
\exists\,\omega\in[\omega_0]\text{ Kähler with constant scalar curvature}
\quad\Longleftrightarrow\quad
L\to X\text{ is stable in the intended geometric invariant theory sense}.
\]
The right-hand side is retained in the source's wording: it does not specify a GIT parameter space, a linearized group action, an asymptotic embedding criterion, or a definition of K-stability. Consequently strict stability, polystability, and their uniform or relative versions cannot be interchanged here. The relationship between \([\omega_0]\) and \(c_1(L)\), including normalization, must also be fixed; the supplied statement does not explicitly give it. No hypothesis on the automorphism group of \((X,L)\) is stated. These choices are necessary to turn the displayed source question into a fully defined algebraic equivalence.
\subsection{Short English statement}
Ask whether a polarized Kähler class contains a constant scalar curvature metric exactly when the polarization satisfies the intended algebraic stability condition. The source leaves the stability notion, automorphism qualifications, and class normalization unresolved.
\subsection{Sources}
\begin{itemize}
\item Local UnsolvedMath registry, permanent ID \(30001003\), source code \texttt{catalogue entry}, ``Yau Stability Conjecture for Constant Scalar Curvature Metrics''; the conjecture is attributed to Yau in the original. Registry curation: Ulam AI / UnsolvedMath Contributors, CC BY 4.0; underlying sources retain their supplied provenance.
\item \emph{Real Analysis, Harmonic Analysis and Applications} (2009), Oberwolfach Reports, pp.\ 1771--1850, DOI \url{https://doi.org/10.4171/owr/2008/32}.
\item Supplied ``Chen--Cheng cscK theorem'': \url{https://arxiv.org/abs/1801.05907}.
\item Supplied ``current K-stability survey'': \url{https://doi.org/10.1017/fms.2023.104}.
\item Supplied literature triage, \texttt{literature\_checked\_at}: 2026-08-21: labeled partially solved; says major cases and analytic directions of the cscK--stability correspondence are proved, while the general algebraic Yau--Tian--Donaldson equivalence in the stated form remains unresolved. This is dated registry curation requiring independent status and hypothesis review.
\end{itemize}
\end{document}
